\section*{Capítulo \ref{ch:b2:ism-star-formation} --- El medio interestelar y la formación de estrellas}

\begin{solution}{exo:b2:ism-star-formation:1}
$n = \num{2.7e-24}/\num{1.67e-24} = 1.6$ átomos de hidrógeno por centímetro cúbico. En el aire,
$n = P/k_{\mathrm B}T = 101325/(\num{1.381e-23}\times288) = \qty{2.55e25}{m^{-3}}$,
es decir \num{2.5e19} moléculas por centímetro cúbico: unas $1.6\times10^{19}$ veces
más.
\end{solution}

\begin{solution}{exo:b2:ism-star-formation:2}
$M_J \propto n^{-1/2}$: dividir $n$ entre 100 multiplica $M_J$ por 10, de modo que
$M_J = 54.6$ masas solares a $10^2$ partículas por centímetro cúbico. La nube contiene
$\num{24000}/54.6 \approx 440$ \omterm{def:b2:ism-star-formation:jeans}{masas de Jeans}: no puede ser estable en su conjunto, y tiene que
fragmentarse.
\end{solution}

\begin{solution}{exo:b2:ism-star-formation:3}
$c_s = \sqrt{\num{1.381e-23}\times10/(2.33\times\num{1.661e-27})} =
\qty{189}{m.s^{-1}}$ a \qty{10}{K}; $\sqrt{10}$ veces más, \qty{597}{m.s^{-1}},
a \qty{100}{K}. Para atravesar \qty{0.1}{pc}: $\num{3.086e15}/189 = \qty{1.63e13}{s}$,
es decir \qty{0.52}{Myr}.
\end{solution}

\begin{solution}{exo:b2:ism-star-formation:4}
$\rho = 2.33\times\num{1.661e-27}\times n$ con $n$ en \unit{m^{-3}}. Para $10^5$
partículas por centímetro cúbico, $\rho = \qty{3.87e-16}{kg.m^{-3}}$ y
$t_{\mathrm{ff}} = \sqrt{3\pi/(32G\rho)} = \qty{3.38e12}{s} = \qty{0.107}{Myr}$; para
$10^2$, $\rho$ es \num{1000} veces menor y $t_{\mathrm{ff}}$ $\sqrt{1000} = 31.6$
veces más largo, \qty{3.38}{Myr}. El \omterm{def:b2:ism-star-formation:molecular-cloud}{núcleo denso} colapsa 32 veces más deprisa.
\end{solution}

\begin{solution}{exo:b2:ism-star-formation:5}
Con $k_{\mathrm B}T/(\mu m_u) = c_s^2$, $(5k_{\mathrm B}T/G\mu m_u)^{3/2} =
5^{3/2}c_s^3G^{-3/2}$, lo que da la fórmula. Como $M_J \propto \rho^{-1/2}$,
reducir $M_J$ a la mitad exige una densidad cuatro veces mayor. Como $M_J \propto T^{3/2}$,
calentar de \qty{10}{K} a \qty{20}{K} la multiplica por $2^{3/2} = 2.83$.
\end{solution}

\begin{solution}{exo:b2:ism-star-formation:6}
$t_{\mathrm{KH}} = GM^2/(RL) = \num{6.674e-11}\times(\num{1.989e30})^2/(2\times
\num{6.96e8}\times2\times\num{3.83e26}) = \qty{2.48e14}{s}$, es decir
\qty{7.8}{Myr}. Es más largo que la duración de la mayoría de los discos (la mitad han desaparecido tras
unos \qty{3}{Myr}): la estrella todavía se está contrayendo cuando su disco desaparece.
\end{solution}

\begin{solution}{exo:b2:ism-star-formation:7}
A \qty{10}{K}, $c_s^3/G = 189^3/\num{6.674e-11} = \qty{1.01e17}{kg.s^{-1}}$, es decir
$\num{1.01e17}\times\num{3.156e7}/\num{1.989e30} = \num{1.6e-6}$ masas solares por
año, y una masa solar tarda \qty{0.62}{Myr}. A \qty{20}{K} la tasa es
$2^{3/2} = 2.83$ veces mayor, \num{4.5e-6} masas solares por año, y una masa
solar tarda \qty{0.22}{Myr}.
\end{solution}

\begin{solution}{exo:b2:ism-star-formation:8}
$r = \sqrt{L/(16\pi\sigma)}/T^2$ con $\sqrt{\num{3.83e26}/(16\pi\times\num{5.67e-8})}
= \qty{1.159e16}{m.K^2}$: a \qty{1300}{K}, $r = \qty{6.86e9}{m} = \qty{0.046}{au}$,
muy dentro de la órbita de Mercurio. A \qty{5.2}{au}, $T = 278/\sqrt{5.2} = \qty{122}{K}$.
\end{solution}

\begin{solution}{exo:b2:ism-star-formation:9}
$v = \sqrt{GM/r} = \sqrt{\num{6.674e-11}\times1.3\times\num{1.989e30}/\num{1.496e13}}
= \qty{3.40}{km.s^{-1}}$, y $P = 2\pi r/v = \qty{2.77e10}{s} = 877$ años
(o, de forma equivalente, $\sqrt{100^3/1.3}$ años). En un millón de años el gas ha dado
$10^6/877 \approx \num{1140}$ vueltas: el disco es viejo comparado con su propia rotación,
y por eso sus anillos son círculos regulares.
\end{solution}

\begin{solution}{exo:b2:ism-star-formation:10}
$M = \int_0^R 4\pi r^2\rho_0R/r\,\mathrm dr = 2\pi\rho_0R^3$, y dentro de $r$,
$m(r) = 2\pi\rho_0Rr^2$. Con $\mathrm dm = 4\pi\rho_0Rr\,\mathrm dr$,
\[
  U = -\int_0^R \frac{G\,m(r)}{r}\,\mathrm dm
    = -8\pi^2G\rho_0^2R^2\int_0^R r^2\,\mathrm dr = -\frac{8\pi^2}{3}G\rho_0^2R^5
    = -\frac23\,\frac{GM^2}{R},
\]
usando $\rho_0 = M/(2\pi R^3)$. Es $\tfrac{2/3}{3/5} = 1.11$ veces el valor uniforme:
la nube concentrada está más ligada y necesita un \qty{11}{\percent} más de
energía térmica para sostenerse; colapsa con más facilidad.
\end{solution}

\begin{solution}{exo:b2:ism-star-formation:11}
Periodo $2\pi/\Omega = \qty{6.28e14}{s}$, unos \qty{20}{Myr}. Con
$R = \qty{3.086e15}{m}$, $j = \Omega R^2 = \qty{9.52e16}{m^2.s^{-1}}$, y
$r_c = j^2/(GM) = \num{9.07e33}/\num{1.327e20} = \qty{6.83e13}{m} = \qty{457}{au}$.
En rotación de sólido rígido $j = \Omega r^2$ crece como el cuadrado de la distancia al eje:
el gas del interior del núcleo tiene menos momento angular, y cae más cerca de la estrella.
\end{solution}

\begin{solution}{exo:b2:ism-star-formation:12}
$\bar t = 2.46$, $\bar f = 0.588$; $\sum(t - \bar t)^2 = 12.612$,
$\sum(t - \bar t)(f - \bar f) = -2.018$, de modo que $b = 2.018/12.612 = 0.160$ por millón
de años y $a = 0.588 + 0.160\times2.46 = 0.982$. La recta llega a $0.5$ en
$(0.982 - 0.5)/0.160 = \qty{3.0}{Myr}$ y a cero en $0.982/0.160 = \qty{6.1}{Myr}$, como en
la figura. Una fracción no puede bajar de cero: los cúmulos más viejos, como NGC 1960 a
\qty{30}{Myr}, conservan una pequeña fracción cercana a cero que una recta no puede describir.
\end{solution}

\begin{solution}{pb:b2:ism-star-formation:1}
\textbf{1.} $R = 0.08\times\num{3.086e16} = \qty{2.469e15}{m}$.

\textbf{2.} $N = \qty{1.5e26}{m^{-2}}$, de modo que $2.8\times\num{1.661e-27}\times\num{1.5e26} =
\qty{0.698}{kg.m^{-2}}$.

\textbf{3.} $M = 0.698\times\pi R^2 = 0.698\times\num{1.915e31} = \qty{1.34e31}{kg}$,
es decir $6.7$ masas solares.

\textbf{4.} $V = \tfrac43\pi R^3 = \qty{6.30e46}{m^3}$ y
$\rho = M/V = \qty{2.12e-16}{kg.m^{-3}}$.

\textbf{5.} $n = \rho/(2.33m_u) = \qty{5.48e10}{m^{-3}}$, es decir \num{5.5e4} partículas
por centímetro cúbico.

\textbf{6.} Sí: su densidad está entre $10^4$ y $10^5$, su radio entre 0.03 y
\qty{0.2}{pc}; su masa, 6.7 masas solares, está en lo alto del intervalo de los núcleos.

\textbf{7.} $M_J = 5.46\times(10^4/\num{5.48e4})^{1/2} = 2.33$ masas solares.

\textbf{8.} $6.7/2.33 = 2.9$ \omterm{def:b2:ism-star-formation:jeans}{masas de Jeans}: el núcleo es inestable y colapsa.

\textbf{9.} $R_J = 0.1316\times(10^4/\num{5.48e4})^{1/2} = \qty{0.056}{pc}$, menor que
el radio, \qty{0.08}{pc}: la misma conclusión.

\textbf{10.} $2^{3/2}\times2.33 = 6.6$ masas solares: aproximadamente la masa del núcleo. A
\qty{20}{K} el núcleo sería marginal; es el frío lo que lo hace colapsar.

\textbf{11.} $t_{\mathrm{ff}} = \sqrt{3\pi/(32\times\num{6.674e-11}\times\num{2.12e-16})}
= \qty{4.56e12}{s}$, unos \num{145000} años.

\textbf{12.} $R/c_s = \num{2.469e15}/189 = \qty{1.31e13}{s}$, \num{414000} años: tres
veces más que la caída libre. Una onda de presión no puede atravesar el núcleo antes de que este haya
colapsado; la presión no puede reaccionar lo bastante deprisa para detener el colapso.

\textbf{13.} $c_s = \qty{189}{m.s^{-1}}$ (\cref{exo:b2:ism-star-formation:3}).

\textbf{14.} $c_s^3/G = \qty{1.01e17}{kg.s^{-1}}$, \num{1.6e-6} masas solares por año.

\textbf{15.} $1/\num{1.6e-6} = \num{624000}$ años.

\textbf{16.} Un tercio de 6.7 es $2.2$: dos estrellas de una masa solar.

\textbf{17.} $t_{\mathrm{KH}} = GM_\odot^2/(3R_\odot\times4L_\odot) = \qty{8.25e13}{s}$,
\qty{2.6}{Myr}.

\textbf{18.} Más largo, cuatro veces el tiempo de acreción: cuando la acreción se detiene, la estrella
todavía se contrae y brilla sin fusión; es una \omterm{def:b2:ism-star-formation:t-tauri}{estrella T Tauri}, que aún no está en la secuencia
principal.

\textbf{19.} $\Omega = \qty{0.4e3}{m.s^{-1}}/\qty{3.086e16}{m} = \qty{1.30e-14}{rad.s^{-1}}$.

\textbf{20.} $j = 1.30\times10^{-14}\times(\num{2.469e15})^2 = \qty{7.90e16}{m^2.s^{-1}}$.

\textbf{21.} $r_c = j^2/(GM_\odot) = \num{6.24e33}/\num{1.327e20} = \qty{4.70e13}{m}$, es
decir \qty{314}{au}: más que el disco de HL Tauri, como cabe esperar para gas del borde del
núcleo.

\textbf{22.} $T = (L/16\pi\sigma r^2)^{1/4} = \qty{278}{K}$ a \qty{1}{au}, y
$278/\sqrt{30} = \qty{51}{K}$ a \qty{30}{au}.

\textbf{23.} $r = \num{1.159e16}/180^2 = \qty{3.58e11}{m}$, es decir \qty{2.4}{au}, como en el texto.

\textbf{24.} La estrella de una masa solar tarda
\textbf{\qty[text-series-to-math=true]{624}{} mil años} en ensamblarse, unas cuatro
veces el \omterm{def:b2:ism-star-formation:free-fall}{tiempo de caída libre} de su núcleo (\num{145000} años): el gas exterior llega a la
\omterm{def:b2:ism-star-formation:protostar}{protoestrella} solo gradualmente, de dentro hacia fuera.
\end{solution}
